This concluding section develops concrete research agendas at the intersection of reinforcement learning and the applied sciences as well as lists the bottlenecks and open challenges.

\subsection{How Domain Structure Improves Reinforcement Learning}

Reinforcement learning's most celebrated successes, from Atari to Go to protein folding, share a common ingredient: a cheap, fast, and accurate simulator that generates unlimited training data. Most applied domains lack this ingredient. Every application in Section~\ref{section:applications} demanded a custom environment, and the engineering cost of building these environments often dominated the cost of training the RL agent itself. Structural models are, however, uniquely positioned to fill this gap. They encode variable selection, causal structure, and institutional constraints that determine how agents respond to interventions. The Lucas critique warns that correlational simulators, trained on observational data without structural assumptions, will break under policy changes (Section~\ref{section:causal_rl}). Building simulators that respect causal identification and correctly model agent responses to rule changes is an important problem.

A recurring theme across the survey is that domain structure, when available, can dramatically reduce the sample complexity of learning. The knowledge ladder in Section~\ref{section:bandits} shows that imposing demand structure on a pricing problem can reduce cumulative regret from $\Theta(T)$ to $O(\log T)$. The pattern extends beyond bandits: structural assumptions yield similar gains in dynamic estimation (Section~\ref{section:rl_econ_models}) and preference learning (Section~\ref{section:offline_rl}). These reflect the difference between learning in an unstructured space and learning in one shaped by theory. Formalizing this intuition, identifying which structural assumptions yield which complexity reductions and under what conditions, is an active research frontier.

The RLHF and DPO frameworks studied in Section~\ref{section:offline_rl} rely on the Bradley-Terry model, one of the simplest members of the discrete choice family. The discrete choice literature offers a rich toolkit for moving beyond it. Mixed logit models accommodate heterogeneous preferences across annotators. Revealed preference theory provides axiomatic consistency constraints, such as GARP and stochastic transitivity, that can serve as regularizers on learned reward models. The preference learning simulation in Section~\ref{section:offline_rl} illustrates the cost of misspecification (Figure~\ref{fig:preference_sample}). Integrating tools from econometrics for preference elicitation, model selection, and specification testing into the RLHF pipeline is a natural direction.

\subsection{How Reinforcement Learning Advances Applied Modeling}

Dynamic programming has always offered a prescriptive capability: given a model, compute the optimal policy. In practice, this capability has been limited by the curse of dimensionality to models with small, discrete state spaces. RL relaxes this constraint. TD-based methods make structural models tractable at state-space scales where NFXP is infeasible (Section~\ref{section:rl_econ_models}), and real-world deployments confirm that RL can compute policies of practical value (Section~\ref{section:applications}). These results suggest a prescriptive role for RL: not merely estimating model parameters (the traditional estimation task), but computing the policies those parameters imply.

Social science and policy research work with vast quantities of observational data from settings where controlled experimentation is impossible or unethical. Offline RL, which learns policies from logged data without further interaction, is a natural fit. The simulation in Section~\ref{section:offline_rl} illustrates both the promise and the limits: algorithms with distributional shift correction exceed the behavioral baseline, while those without it degrade below it, and performance depends critically on data support (Table~\ref{tab:offline_main}, Figure~\ref{fig:offline_coverage}). This means that offline RL is not a black box that extracts optimal policies from any observational dataset. It is a framework with clearly delineated conditions under which learning is possible, conditions that connect directly to familiar statistical concepts like overlap and common support. Developing standardized benchmarks and digital twins for offline RL can enable firms and policy-makers to design optimal policies from data generated under old ones.

RL also offers a descriptive model of how boundedly rational agents learn in strategic environments. The simulations in Section~\ref{section:rl_games} demonstrate that independent Q-learning agents converge to Nash equilibrium through trial and error, providing ``as-if'' microfoundations for equilibrium concepts: the equilibrium arises not from common knowledge of rationality, but from a simple adaptive process.

When RL-trained agents are actually deployed in markets, they become market actors subject to empirical scrutiny. Algorithmic pricing agents, for instance, have been shown to sustain supra-competitive prices through reward-based learning without explicit communication. This is not a speculative concern; competition authorities in the EU, US, and UK are actively investigating whether algorithmic coordination constitutes tacit collusion.\footnote{The companion thesis \citep{Rawat2026collusion} develops a framework for evaluating algorithmic inefficiency and collusion risk in algorithmically mediated markets, combining simulators with factorial experimental designs.} As algorithmic agents proliferate in pricing, trading, content recommendation, and resource allocation, the empirical study of machine behavior becomes an increasingly important research program. 

\subsection{Open Challenges}

For researchers accustomed to estimators with well-characterized asymptotic properties, the deadly triad represents a significant barrier to adoption. Deep RL algorithms exhibit seed sensitivity, overestimation cascades, and plasticity loss that make reproducibility difficult even in controlled settings (Section~\ref{section:deeprl_practice}).

Multi-agent RL faces fundamental problems (Section~\ref{section:rl_games}). Computing Nash equilibria is PPAD-complete, and RL does not escape this hardness. In games with multiple equilibria, agents performing Bellman backups may select different equilibria for different states, causing value iteration to cycle or diverge. The literature on equilibrium refinement, focal points, and coordination mechanisms may offer partial resolutions, but a general solution remains open.

The absence of standardized simulators is part of a broader infrastructure gap. Industrial RL deployments require dedicated engineering teams, GPU clusters, and months of hyperparameter tuning (Section~\ref{section:applications}, Section~\ref{section:deeprl_practice}). Reducing these barriers through shared simulation environments and accessible software libraries would accelerate adoption.

\subsection{Conclusion}

Reinforcement learning is a welcome addition to the researcher's toolkit but is not a replacement for existing methods. It extends the reach of dynamic programming to problems that were previously intractable, and it connects naturally to established statistical frameworks through shared mathematical foundations. The research agendas outlined above, from building structural simulators to formalizing the role of structural assumptions in sample complexity to developing inference procedures for learned policies, are concrete and tractable. Progress on them will require sustained collaboration between the two fields, drawing on domain knowledge for structure and institutional context and on RL for computation.
